%==============================================================
%  CAPITOLO 1 - PREFAZIONE
%==============================================================
\chapter{Prefazione}
\label{cap:prefazione}

\textbf{Sicurezza funzionale dei comandi delle macchine}\\
\textbf{Applicazione della norma EN ISO 13849-1}

La norma EN ISO 13849-1, "Sicurezza delle macchine -Parti dei sistemi di controllo relative alla sicurezza", contiene disposizioni che regolano la progettazione di tali parti. Il presente rapporto è un aggiornamento del rapporto BGIA 2/2008e omonimo. Descrive l'argomento essenziale della norma nella terza edizione riveduta del 2015 e ne spiega l'applicazione con riferimento a numerosi esempi tratti dai settori dell'elettromeccanica, della fluidica, dell'elettronica e dell'elettronica programmabile, compresi i sistemi di controllo che impiegano tecnologie miste. La norma viene inserita nel contesto dei requisiti essenziali di sicurezza della Direttiva Macchine e vengono presentati possibili metodi di valutazione del rischio.

Sulla base di queste informazioni, il rapporto può essere utilizzato per selezionare il livello di prestazione PL\textsubscript{\small r} richiesto per le funzioni di sicurezza nei sistemi di controllo. Il livello di prestazione PL effettivamente raggiunto viene spiegato in dettaglio. I requisiti per il raggiungimento del livello di prestazione pertinente e delle relative categorie, l'affidabilità dei componenti, i livelli di copertura diagnostica, la sicurezza del software e le misure per la prevenzione di guasti sistematici e a causa comune vengono discussi in modo esaustivo. Vengono inoltre fornite informazioni di base sull'implementazione dei requisiti nei sistemi di controllo reali. Numerosi circuiti di esempio mostrano, fino al livello dei componenti, come i livelli di prestazione da a ad e possono essere realizzati nelle tecnologie selezionate con le categorie da B a 4. Gli esempi forniscono informazioni sui principi di sicurezza impiegati e sui componenti con funzionalità di sicurezza collaudate. Numerosi riferimenti bibliografici consentono uno studio più approfondito degli esempi forniti. Il rapporto mostra come i requisiti della norma EN ISO 13849-1 possano essere implementati nella pratica ingegneristica, contribuendo così a un'applicazione e un'interpretazione coerenti della norma a livello nazionale e internazionale.
